\section*{Capítulo \ref{ch:b2:plant-meristems} --- Desarrollo vegetativo de las plantas: meristemos y crecimiento}

\begin{solution}{exo:b2:plant-meristems:1}
Un \omterm{def:b2:plant-meristems:meristem}{meristemo} es una población persistente de células no diferenciadas en
división de la que derivan los tejidos de la planta. \omterm{def:b2:plant-meristems:meristem}{Meristemo apical
caulinar}: hojas, tallo, \omterm{def:b2:angiosperm-reproduction:flower}{flores}; \omterm{def:b2:plant-meristems:meristem}{meristemo apical radicular}: la raíz;
\omterm{def:b2:plant-meristems:wood}{cámbium vascular}: xilema secundario (\omterm{def:b2:plant-meristems:wood}{madera}) y floema; felógeno: la
\omterm{def:b2:plant-meristems:wood}{peridermis} (corcho) de la \omterm{def:b2:plant-meristems:wood}{corteza}.
\end{solution}

\begin{solution}{exo:b2:plant-meristems:2}
Cofia (protección, lubricación, percepción de la gravedad, células que se
desprenden); \omterm{def:b2:plant-meristems:meristem}{meristemo} con el \omterm{prop:b2:plant-meristems:ram}{centro quiescente} (división); \omterm{prop:b2:plant-meristems:ram}{zona de
elongación} (las células se alargan de diez a veinte veces); zona de
\omterm{def:b2:cell-differentiation:differentiation}{diferenciación} (\omterm{prop:b2:plant-meristems:ram}{pelos radiculares}, xilema y floema maduros). Las raíces
laterales surgen del periciclo de la raíz madura, frente a los polos del
xilema.
\end{solution}

\begin{solution}{exo:b2:plant-meristems:3}
Se decapitaron plantas de judía: las yemas laterales brotaron. En una
segunda serie, el muñón cortado recibió un bloque de agar con auxina: las
yemas permanecieron latentes; en la serie de control, un bloque de agar
sin auxina: las yemas brotaron. Conclusión: la auxina procedente del ápice
es lo que inhibe las yemas laterales --- la \omterm{prop:b2:plant-meristems:apical}{dominancia apical} es hormonal.
\end{solution}

\begin{solution}{exo:b2:plant-meristems:4}
Sesenta años; la albura son los quince últimos anillos; están vivos el
\omterm{def:b2:plant-meristems:meristem}{cámbium}, el floema, las células de los radios y del parénquima de la
albura y el felógeno --- una fina envoltura alrededor de un núcleo muerto.
\end{solution}

\begin{solution}{exo:b2:plant-meristems:5}
$0.1\times(P - 0.3)$: \qty{0.01}{h^{-1}} (duplicación en $\ln 2/0.01 =
\qty{69}{h}$), \qty{0.03}{h^{-1}} (\qty{23}{h}), \qty{0.05}{h^{-1}}
(\qty{14}{h}). A \qty{0.25}{MPa} la turgencia está por debajo del umbral
de fluencia: el crecimiento se detiene, antes de cualquier marchitamiento.
\end{solution}

\begin{solution}{exo:b2:plant-meristems:6}
Hojas 1--8 a 0, 137.5, 275, 52.5, 190, 327.5, 105, \qty{242.5}{\degree}.
Ordenadas: 0, 52.5, 105, 137.5, 190, 242.5, 275, 327.5 --- el hueco más
pequeño es de \qty{32.5}{\degree}. Cada hoja se sitúa en un hueco de las
anteriores, de modo que las ocho se sombrean lo menos posible y la copa
llena de forma uniforme el disco de luz.
\end{solution}

\begin{solution}{exo:b2:plant-meristems:7}
$2\pi r h\Delta r$: con \qty{5}{cm}, $2\pi\times 0.05\times 12\times
0.003 = \qty{0.011}{m^3}$, \qty{5.7}{kg}, \qty{2.8}{kg} de carbono; con
\qty{40}{cm}, \qty{0.090}{m^3}, \qty{45}{kg}, \qty{23}{kg} de carbono
--- ocho veces más con la misma anchura de anillo.
\end{solution}

\begin{solution}{exo:b2:plant-meristems:8}
\emph{clv3}: sin freno --- \omterm{def:b2:plant-meristems:meristem}{meristemos} enormes, hojas y órganos florales
adicionales, tallos fasciados. \emph{wus}: sin acelerador --- el \omterm{def:b2:plant-meristems:meristem}{meristemo}
se agota tras unas pocas hojas y se refunda una y otra vez. Doble
mutante: el \omterm{def:b2:meiosis-heredity:vocabulary}{fenotipo} \emph{wus}, de modo que WUS actúa aguas abajo de
CLV3: la función de CLV3 es reprimir WUS, y sin WUS no hay nada que
reprimir.
\end{solution}

\begin{solution}{exo:b2:plant-meristems:9}
Transporte: \qty{20}{h}. Difusión: $(0.2)^{2}/(2\times 10^{-9}) =
\qty{2e7}{s}$, unos 230 días. El transporte polar hace utilizable en menos
de un día una señal hormonal a lo largo de una planta, y le da una
dirección.
\end{solution}

\begin{solution}{exo:b2:plant-meristems:10}
\qty{3}{cm} al día de tejido maduro con células de \qty{200}{\micro m}:
\num{150} células por fila al día. Cien células en división deben, por
tanto, dividirse cada una 1,5 veces al día: un ciclo de unas \qty{16}{h}.
\end{solution}

\begin{solution}{exo:b2:plant-meristems:11}
La giberelina alarga los entrenudos; una planta que no puede responder a
ella tiene entrenudos cortos y un tallo bajo. Un abonado nitrogenado
intenso hace crecer un trigo alto más alto y con espigas más pesadas,
hasta que se tumba (se encama) y se pudre; un tallo corto y rígido se
mantiene en pie y destina el carbono adicional al grano en lugar de a la
paja, de modo que el índice de cosecha aumenta. En la naturaleza, una
planta baja queda por debajo de sus vecinas, que la sombrean, y pierde.
\end{solution}

\begin{solution}{exo:b2:plant-meristems:12}
Los \omterm{def:b2:plant-meristems:meristem}{meristemos} permiten a una planta añadir órganos donde y cuando las
condiciones los favorecen --- hojas hacia la luz, raíces hacia el agua ---
y abandonar los que no compensan; cada yema axilar es un punto de
crecimiento de repuesto, de modo que sobrevive al pastoreo, al fuego y a
la poda rebrotando; y como los \omterm{def:b2:plant-meristems:meristem}{meristemos} son perpetuamente embrionarios,
un \omterm{def:b2:asexual-reproduction:clone}{geneto} no tiene una duración de vida fija. El coste: una planta no
puede moverse ni reorganizar lo que ha construido --- un órgano, una vez
colocado, se queda donde está, y un cuerpo construido por acumulación
debe cargar con su pasado muerto (el duramen).
\end{solution}

\begin{solution}{pb:b2:plant-meristems:1}
\textbf{1.} $150/3 = 50$ hojas por ápice; \num{10000} para el árbol.
\textbf{2.} 0, 137.5, 275, 52.5 y \qty{190}{\degree}.
\textbf{3.} $8\times 137.5 = 1100 = 3\times 360 + 20$; $13\times 137.5
= 1787.5 = 5\times 360 - 12.5$; $21\times 137.5 = 2887.5 = 8\times 360
+ 7.5$: la hoja 22 es la primera a menos de \qty{10}{\degree} de la
primera.
\textbf{4.} Cada hoja nueva se sitúa en el hueco más amplio, de modo que
las hojas se sombrean lo menos posible; es la auxina, cuyo transporte
polar (transportadores PIN) drena el entorno de cada primordio, la que
fija el espaciado.
\textbf{5.} $\num{10000}\times 80\,\text{cm}^{2} = \qty{80}{m^2}$: toda la
copa está formada por las hojas de esta estación.
\textbf{6.} Toda; las primeras hojas se construyen con el almidón
almacenado en la \omterm{def:b2:plant-meristems:wood}{madera} y en las raíces el año anterior.
\textbf{7.} Un ápice triplicado forma primordios más deprisa --- dos o tres
por plastocrono ---, de modo que las hojas y la superficie foliar podrían
triplicarse, pero los tallos serían fasciados, con tallos acintados y una
filotaxis desordenada, y buena parte de la hoja adicional se sombrearía a
sí misma.
\textbf{8.} Día: $0.05\times(0.55 - 0.35) = \qty{0.01}{h^{-1}}$; noche:
$0.05\times 0.40 = \qty{0.02}{h^{-1}}$.
\textbf{9.} Día: $2e^{0.12} = \qty{2.25}{cm}$, \qty{0.25}{cm} añadidos;
noche: $2.25e^{0.24} = \qty{2.87}{cm}$, \qty{0.61}{cm} añadidos.
\textbf{10.} Tasa media de \qty{0.015}{h^{-1}} durante \qty{96}{h}: $2e^{1.44}
= \qty{8.4}{cm}$.
\textbf{11.} Día: $0.05\times 0.30 = \qty{0.015}{h^{-1}}$; noche:
$0.05\times 0.50 = \qty{0.025}{h^{-1}}$.
\textbf{12.} $0.05\times 0.05 = \qty{0.0025}{h^{-1}}$, la cuarta parte de la
tasa diurna normal. El ajuste osmótico eleva el contenido de solutos de la
célula de modo que, con el mismo potencial hídrico, la turgencia vuelve a
subir por encima del umbral.
\textbf{13.} De día la transpiración reduce el potencial hídrico de la hoja
y, con él, la turgencia, de modo que $P - Y$ es pequeño; de noche la
turgencia se restablece. Hora a hora, el crecimiento está limitado por el
agua, no por el azúcar.
\textbf{14.} $2\pi\times 0.10\times 8\times 0.005 = \qty{0.025}{m^3}$;
\qty{12.6}{kg} en seco.
\textbf{15.} \qty{6.3}{kg} de carbono.
\textbf{16.} $80\times 4\times 150 = \qty{48}{kg}$ de carbono.
\textbf{17.} $6.3/48 = \qty{13}{\%}$. Otros sumideros: las propias hojas,
las raíces, las ramas y ramillas, la respiración de todas las células
vivas (alrededor de la mitad de lo fijado), las reservas, las \omterm{def:b2:angiosperm-reproduction:flower}{flores} y las
\omterm{def:b2:angiosperm-reproduction:seed}{semillas}.
\textbf{18.} Radio de \qty{0.20}{m}: $2\pi\times 0.20\times 8\times
0.005 = \qty{0.050}{m^3}$, el doble que este año --- la circunferencia a lo
largo de la cual trabaja el \omterm{def:b2:plant-meristems:meristem}{cámbium} se ha duplicado (y el árbol será más
alto).
\textbf{19.} Un año seco o frío; o una defoliación por insectos, heladas o
fuego. Una causa climática estrecha los mismos anillos en todos los
árboles de la región y se correlaciona con los registros meteorológicos;
un daño afecta a un árbol o a un rodal, y las heladas o el fuego dejan
cicatrices anatómicas en el anillo.
\textbf{20.} Las yemas más cercanas a cada corte brotan en pocos días y dan
varios brotes; al cabo de un mes el árbol es más frondoso y denso, con más
brotes, pero más cortos.
\textbf{21.} La auxina en el corte sustituye a la señal del ápice: las
yemas permanecen latentes y no se produce ramificación.
\textbf{22.} Cada recorte elimina los ápices y libera las yemas situadas
debajo; seis veces al año, cada liberación añade un piso de brotes cortos,
y la superficie se llena de ellos.
\textbf{23.} De yemas axilares latentes del tocón y de yemas adventicias
formadas por el \omterm{def:b2:plant-meristems:meristem}{cámbium}; cada yema es un \omterm{def:b2:plant-meristems:meristem}{meristemo} completo capaz de
reconstruir el sistema caulinar, mientras que la cabeza de un animal es
un \omterm{def:b2:vertebrate-development:induction}{organizador} único e insustituible, sin copias de reserva.
\textbf{24.} Más citoquinina llega a las yemas y favorece su brote ---
regar una planta podada la hace ramificar más.
\textbf{25.} 50 hojas por ápice, \num{10000} por árbol; tasas del entrenudo
de \qty{0.01}{h^{-1}} de día y \qty{0.02}{h^{-1}} de noche; la \omterm{def:b2:plant-meristems:wood}{madera} del
año, \qty{0.025}{m^3}, \qty{12.6}{kg} en seco, \qty{6.3}{kg} de carbono.
\end{solution}
